\usepackage{microtype}
\usepackage{fancyhdr}
\usepackage{etoolbox}
\usepackage{xurl}
\usepackage{seqsplit}
\let\oTexttt\texttt
\renewrobustcmd{\texttt}[1]{\oTexttt{\seqsplit{#1}}}
\setlength{\emergencystretch}{3em}
\raggedbottom
\pagestyle{fancy}
\fancyhf{}
\fancyhead[L]{\small Kerja Matematika yang Dapat Ditelusuri}
\fancyhead[R]{\small O017 / D120}
\fancyfoot[C]{\thepage}
\renewcommand{\headrulewidth}{0.3pt}
\renewcommand{\footrulewidth}{0pt}
\AtBeginEnvironment{longtable}{\small}
\newenvironment{refproof}{\begin{proof}}{\end{proof}}
\publishers{%
  \begin{minipage}{0.88\textwidth}
  \centering\small
  \textbf{DOI versi:} \href{https://doi.org/10.5281/zenodo.22073823}{10.5281/zenodo.22073823}\\
  \textbf{DOI semua versi:} \href{https://doi.org/10.5281/zenodo.22051866}{10.5281/zenodo.22051866}\\[0.5em]
  \textbf{Pembaca HTML aksesibel:} \url{https://kokunoyumeto.github.io/kerja-matematika-yang-dapat-ditelusuri-id/}\\
  \textbf{Sumber GitHub:} \url{https://github.com/KokunoYumeto/kerja-matematika-yang-dapat-ditelusuri-id}\\
  \textbf{Arsip edisi, sumber, dan backend:} \url{https://zenodo.org/records/22073823}\\
  \textbf{Batas aksesibilitas:} PDF ini tidak bertag dan bukan edisi aksesibel; gunakan pembaca HTML.\\[0.7em]
  OpenAI Codex gpt-5.6-sol, Ultra. Digunakan atas instruksi kreator manusia.\\
  Pemetaan hak per komponen tersedia dalam \texttt{LICENSES.md}.
  \end{minipage}%
}
